\chapter{두 번째 출력}
% full version: chapters/16_chemoutput.tex

기계에는 출력이 둘 있다. 빠른 출력은 근육이다. 느린 출력은 몸의 화학이다. 뇌는 심장, 혈관, 샘을 제어하고, 혈액을 통해 온몸에 한꺼번에 메시지를 보낸다. 혈액은 모든 버스 가운데 가장 넓게 퍼지는 브로드캐스트 버스다. 메시지 하나가 약 1분 만에 모든 세포에 닿고, 그 의미는 언제나처럼 수신자의 “자물쇠”가 정한다.

\section*{가속과 브레이크}

심장은 스스로 뛴다. 심장에는 자기 메트로놈이 있다. 뇌는 부호가 반대인 전선 두 개로 박자를 조금씩 조정할 뿐이다. \nt{NORA}는 가속 페달이다. 박자를 빠르게 하지만 느리게 작동한다. \nt{ACET}은 브레이크다. 긴 화학 반응 사슬이 필요 없어서 빠르게 박자를 늦춘다. 그리고 여기서 마침내 \nt{ADRE} 유형이 할 일이 나온다. 혈액으로 뿜어진 아드레날린은 같은 가속 페달이되, 온몸에 대해 한꺼번에 밟은 페달이다.

\pencilfig{16}{\pencilart{16}}{심장은 부호가 반대인 두 전선이 조종한다. 가속과 브레이크다.}

\section*{되먹임이 있는 캐스케이드}

많은 호르몬은 캐스케이드로 제어된다. 뇌가 첫 번째 샘에 명령하고, 그 샘이 두 번째 샘에 명령하고, 두 번째 샘이 호르몬을 분비하면, 호르몬은 이제 충분하다고 뇌에 알린다. 3단 온도 조절기다. 엔지니어는 이런 회로의 위험을 안다. 루프를 너무 “민감하게” 만들면 흔들리기 시작한다. 수준이 유지되지 않고 진동한다. 그래서 이런 조절기의 정확도에는 한계가 있다. 안정하면서 동시에 아주 정확하게 만들 수는 없다.

\pencilfig{127}{%
% three identical first-order stages with proportional feedback: secretion = max(0, K (target - level)).
% K = 3 settles at 3/4 of the target; K = 12 (above the stability limit K = 8) keeps oscillating. x = 0.42 t, y = 0.55 level
\foreach \yo/\t in {1.75/{약한 루프: 안정하지만 목표에 못 미침}, 0/{강한 루프: 목표에 가깝지만 흔들림}}{
  \draw[pen] (0,\yo) -- (8.5,\yo);
  \draw[pcGray, densely dashed, line width=0.5pt] (0,\yo+0.55) -- (8.5,\yo+0.55);
  \node[lbl, anchor=south west] at (2.0,\yo+0.8) {\t};}
\node[lbl, anchor=west] at (8.55,2.3) {목표}; \node[lbl, anchor=west] at (8.55,0.55) {목표};
\begin{scope}[yshift=1.75cm]
\draw[penlite=pcBlue] plot[smooth] coordinates {(0.00,0.00) (0.17,0.01) (0.34,0.08) (0.50,0.19) (0.67,0.33) (0.84,0.46) (1.01,0.55) (1.18,0.59) (1.34,0.58) (1.51,0.55) (1.68,0.49) (1.85,0.43) (2.02,0.37) (2.18,0.34) (2.35,0.33) (2.52,0.34) (2.69,0.37) (2.86,0.40) (3.02,0.43) (3.19,0.44) (3.36,0.45) (3.53,0.45) (3.70,0.44) (3.86,0.42) (4.03,0.41) (4.20,0.40) (4.37,0.39) (4.54,0.39) (4.70,0.40) (4.87,0.40) (5.04,0.41) (5.38,0.42) (5.88,0.42) (6.38,0.41) (7.06,0.41) (8.40,0.41)};
\end{scope}
\draw[penlite=pcRed] plot[smooth] coordinates {(0.00,0.00) (0.17,0.05) (0.34,0.30) (0.50,0.68) (0.67,1.02) (0.84,1.23) (1.01,1.30) (1.18,1.26) (1.34,1.16) (1.51,1.02) (1.68,0.87) (1.85,0.72) (2.02,0.58) (2.18,0.47) (2.35,0.39) (2.52,0.38) (2.69,0.46) (2.86,0.57) (3.02,0.67) (3.19,0.71) (3.36,0.70) (3.53,0.65) (3.70,0.58) (3.86,0.50) (4.03,0.43) (4.20,0.41) (4.37,0.45) (4.54,0.53) (4.70,0.61) (4.87,0.65) (5.04,0.64) (5.21,0.60) (5.38,0.54) (5.54,0.47) (5.71,0.42) (5.88,0.42) (6.05,0.48) (6.22,0.56) (6.38,0.62) (6.55,0.64) (6.72,0.62) (6.89,0.56) (7.06,0.50) (7.22,0.44) (7.39,0.42) (7.56,0.45) (7.73,0.53) (7.90,0.60) (8.06,0.63) (8.23,0.62) (8.40,0.58)};
\node[lbl, anchor=north] at (4.2,-0.05) {시간};
}{3단 호르몬 캐스케이드의 모델. 약한 되먹임은 수준을 안정시키지만 목표보다 눈에 띄게 낮게 남긴다. 강한 되먹임은 목표에 더 가깝게 데려가는 대신, 수준이 머물지 못하고 흔들린다.}

\section*{수준이 아니라 펄스}

어떤 호르몬은 펄스로만 작동한다. 계속 흘려보내면 수신 쪽 샘이 “지쳐서” 꺼져 버리고, 한 시간에 한 번씩 짧은 펄스로 보내면 작동한다. 수신자는 수준이 아니라 리듬을 읽는다. 모스 부호를 다룬 장에서 본, 변화만 듣는 시냅스와 같다.

\section*{하루}

마지막으로, 몸에는 주기가 약 하루인 느린 발진기가 있다. 이 발진기의 고유 주기는 24시간과 조금 다르고, 매일 아침 빛이 이것을 맞춰 준다. 라디오 수신기를 방송국에 맞추는 것과 같다. 이것을 컴퓨터의 클록 발생기와 혼동하지 말자. 이 발진기는 계산의 단계를 재지 않는다. 낮과 밤, 깨어 있음과 잠처럼 기계 전체의 동작 모드를 전환한다.

\fullversion{16}
